\documentclass[10pt,a4paper]{article}

% Language and Encoding
\usepackage[utf8]{inputenc}
\usepackage[T1]{fontenc}
\usepackage[ngerman]{babel}

% Font (Sans Serif to match the modern CV look)
\usepackage{helvet}
\renewcommand{\familydefault}{\sfdefault}

% Layout and Margins (Reduced outer borders so more content fits cleanly)
\usepackage[top=0.8cm, bottom=0.8cm, left=0.9cm, right=0.9cm]{geometry}
\usepackage{graphicx}
\usepackage{multicol}
\usepackage{fontawesome5}
\usepackage{enumitem}
\usepackage{hyperref}
\usepackage{titlesec}
\usepackage{xcolor}

% Custom Colors
\definecolor{dark}{RGB}{0,0,0}

% Section formatting
\titleformat{\section}
  {\large\bfseries\color{dark}}
  {}{0em}
  {\uppercase}[\vspace{-0.5em}\rule{\linewidth}{0.8pt}]
\titlespacing{\section}{0pt}{1.2ex}{0.8ex}

% Custom list spacing for tighter bullets
\setlist[itemize]{leftmargin=1.4em, itemsep=1.5pt, parsep=0pt, topsep=1.5pt}

% Remove page numbers
\pagestyle{empty}

\begin{document}

% ==========================================
% HEADER SECTION (Picture + Profile Summary)
% ==========================================
\begin{minipage}[c]{0.14\textwidth}
    \includegraphics[width=\linewidth]{profile.jpg}
\end{minipage}
\hfill
\begin{minipage}[c]{0.80\textwidth}
    {\Huge \textbf{Iheb Elleuch}}\\[0.3em]
    {\large \textbf{\textit{B.Eng. Automatisierungssysteme und Produktionsinformatik}}}\\[0.5em]
    \hspace*{1.5em}Absolvent der Automatisierungstechnik, spezialisiert auf industrielle und hardwarenahe Softwareentwicklung, Visualisierung und Überwachung. Erforschte LLM-basierte Lösungen zur Erstellung von Prozessdateien in der Laserfertigung mit begrenzten Daten.\\[0.4em]
    \hspace*{1.5em}Nachgewiesene Fähigkeit, Projekte zu strukturieren, eigenverantwortlich zu arbeiten und Ergebnisse termingerecht zu liefern. Aktueller Masterstudent der Technischen Kybernetik an der Universität Stuttgart mit Schwerpunkt auf Maschinelles Lernen, Robotik und Regelungstechnik.
\end{minipage}

\vspace{0.7em}

% ==========================================
% CONTACT INFORMATION
% ==========================================
\hrule height 0.8pt
\vspace{0.1em}
\begin{center}
    \faEnvelope \ \href{mailto:elleuchiheb2002@gmail.com}{elleuchiheb2002@gmail.com} \quad
    \faPhone \ +49 176 84848534 \quad
    \faLinkedin \ \href{https://de.linkedin.com/in/iheb-elleuch}{linkedin.com/in/iheb-elleuch} \quad
    \faGlobe \ \href{https://elleuch-tech.de/}{elleuch-tech.de} \quad
    \faMapMarker \ 73734 
\end{center}
\vspace{0.1em}
\hrule height 0.4pt
\vspace{1.0em}

% ==========================================
% BODY SECTION (Two Columns)
% ==========================================
\begin{minipage}[t]{0.66\textwidth}
    
    % --- EDUCATION ---
    \section*{AKADEMISCHER WERDEGANG}
    \textbf{Master of Science} \hfill seit 10/2026\\
    Universität Stuttgart (Stuttgart-Vaihingen, Deutschland)\\[0.2em]
    \textit{Technische Kybernetik}
    \begin{itemize}
        \item \textbf{Schwerpunkte:} Konzepte der Regelungstechnik (Lyapunov-Stabilität, LQR, $\mathcal{H}_{\infty}$), Modellierung \& Simulation in der Mechatronik, Maschinelles Lernen in der Systemdynamik, Regelungstechnisches Praktikum (3DOF Helikopter).
    \end{itemize}
    \vspace{0.6em}

    \textbf{Bachelor of Engineering} \hfill 09/2022 - 02/2026\\
    Hochschule Esslingen (Göppingen, Deutschland)\\[0.2em]
    \textit{Automatisierungssysteme und Produktionsinformatik}
    \begin{itemize}
        \item \textbf{Abschlussnote:} 1,9 (210 ECTS)
        \item \textbf{Kernqualifikationen:} Echtzeit-Softwareentwicklung, verteilte Automatisierungslösungen und industrielle Netzwerke (SPS, Feldbusse, Motion Control und Ferndiagnose).
        \item \textbf{Ingenieurwissenschaftlicher Fokus:} Hardware/Software-Codesign, Simulation und Inbetriebnahme mechatronischer Systeme durch interdisziplinäre Projektierung.
        \item \textbf{Auslandssemester:} ESTACA, Paris, Frankreich (01/2025 - 05/2025)
        \item \textbf{Engagement:} Mitgründer der Unterwasser-Roboter-Initiative U.Stall
    \end{itemize}
    
    \vspace{0.8em}

    % --- WORK EXPERIENCE ---
    \section*{BERUFSERFAHRUNG}
    
    \textbf{Bachelor-Thesis im Bereich Lasertechnik} \hfill 09/2025 - 02/2026\\
    Bosch Manufacturing Solutions GmbH
    \vspace{-0.2em}
    \begin{itemize}
        \item Nachweis der Machbarkeit von LLMs zur Unterstützung bei der Erstellung von Prozessdateien im industriellen Umfeld.
        \item Überwindung begrenzter Datenverfügbarkeit durch Techniken zur synthetischen Datenskalierung und Entwicklung einer Benchmark-Methode.
        \item Evaluierung von Fine-Tuning-Techniken und Daten-Kurations-Methoden zur Leistungsoptimierung auf lokaler Hardware.
    \end{itemize}
    
    \vspace{0.6em}
    \textbf{Praktikant in Research \& Software Engineering} \hfill 09/2024 - 01/2025\\
    Bosch Manufacturing Solutions GmbH
    \vspace{-0.2em}
    \begin{itemize}
        \item Leitung der UI/UX-Entwicklung einer C\#/WPF-Benutzeroberfläche für die dynamische Laserstrahlformung.
        \item Entwurf einer robusten Backend-Kommunikationsarchitektur mit gRPC und Python.
        \item Erweiterung der Backend-Funktionalität zur Unterstützung neuer Features und Optimierung von Laserstrahlformungs-Algorithmen.
        \item Durchführung routinemäßiger Validierungen des Software-Prototyps.
    \end{itemize}

    \vspace{0.6em}
    \textbf{Werkstudent - Echtzeitsysteme} \hfill 12/2022 - 08/2024\\
    Schuler Pressen GmbH (Andritz)
    \vspace{-0.2em}
    \begin{itemize}
        \item Erweiterung der Überwachungsinfrastruktur und Visualisierung kritischer Daten durch Grafana-Dashboards.
        \item Erfassung und Optimierung von Datenpipelines mittels Prometheus, PostgreSQL und Node-RED-Routing.
        \item Entwicklung und Wartung zeitsparender Automatisierungsskripte für schnellere Datenverarbeitung und vereinfachte Upgrade-Iterationen.
    \end{itemize}

\end{minipage}
\hfill
\begin{minipage}[t]{0.30\textwidth}

    % --- SKILLS ---
    \section*{KOMPETENZEN}
    \begin{itemize}
        \item \textbf{Programmierung:} C++, C\#, Python, JavaScript, SQL, MATLAB/Simulink, VBA
        \item \textbf{Technologien:} WPF, .NET, gRPC, Linux, Node-RED, Grafana, Prometheus, PostgreSQL, LLM Fine-Tuning
        \item \textbf{Robotik \& Regelung:} Hardware-in-the-Loop, Sensor-/Aktor-Integration, SPS, Motion Control, Echtzeitsysteme
        \item \textbf{Werkzeuge:} Git (GitHub, GitLab, Azure), Docker, Visual Studio, VS Code, Unity Engine, Axolotl, Llama-Factory
        \item \textbf{Staatsangehörigkeit:} Deutsch, Tunesisch
        \item \textbf{Sonstiges:} Führerschein Klasse B
    \end{itemize}
    
    % --- LANGUAGES ---
    \section*{SPRACHEN}
    \begin{itemize}
        \item \textbf{Deutsch:} Muttersprache
        \item \textbf{Arabisch:} Muttersprache
        \item \textbf{Englisch:} C1 Verhandlungssicher
        \item \textbf{Französisch:} B2 Fortgeschritten
    \end{itemize}
    
    \vspace{0.6em}
    
    % --- PROJECTS ---
    \section*{SCHLÜSSELPROJEKTE}
    \begin{itemize}
        \item \textbf{U.Stall (MATE ROV):}\\
        Software Lead \& Team Lead (Eine Saison). Architektur und Entwicklung des kompletten Software-Stacks für die vertikale Integration zwischen dem ROV und der Piloten-Schnittstelle. Internationale Teilnahme: USA 2024 (Pioneer), USA 2025 (Explorer), Kanada 2026 (Explorer).
        \vspace{0.4em}
        \item \textbf{Akademische Robotikprojekte:}\\
        \textbf{PinguBlocks:} Mitentwicklung eines sprachgesteuerten, pinguinförmigen Roboters (Raspberry Pi, Webinterface).\\
        \textbf{KlemmbausteinIoT:} Mitentwicklung einer modularen IoT-Robotikplattform aus Klemmbausteinen und handelsüblicher Hardware für den akademischen Unterricht.
    \end{itemize}

\end{minipage}

\end{document}